% Shared preamble for the robodog engineering documentation.
\usepackage[a4paper,margin=24mm,top=22mm,bottom=24mm]{geometry}
\usepackage[T1]{fontenc}
\usepackage[utf8]{inputenc}
\usepackage{lmodern}
\usepackage{microtype}
\usepackage{graphicx}
\usepackage{booktabs}
\usepackage{longtable}
\usepackage{array}
\usepackage{amsmath}
\usepackage{xcolor}
\usepackage{listings}
\usepackage{tcolorbox}
\usepackage{enumitem}
\usepackage{caption}
\usepackage{float}
\usepackage[hidelinks,bookmarks=true]{hyperref}

\definecolor{rdblue}{HTML}{1F4E79}
\definecolor{rdgrey}{HTML}{5A6675}
\definecolor{rdbg}{HTML}{F4F6F9}
\definecolor{rdline}{HTML}{C8D0DA}
\definecolor{rdred}{HTML}{A4262C}

\hypersetup{linkcolor=rdblue, citecolor=rdblue, urlcolor=rdblue,
            colorlinks=true, pdftitle={robodog -- software architecture}}

\captionsetup{font=small, labelfont={bf,color=rdblue}, width=0.92\textwidth}
\setlist{nosep, leftmargin=1.4em}
\renewcommand{\arraystretch}{1.15}

\lstset{
  basicstyle=\ttfamily\footnotesize,
  backgroundcolor=\color{rdbg},
  frame=single, rulecolor=\color{rdline},
  breaklines=true, showstringspaces=false,
  keywordstyle=\color{rdblue}\bfseries,
  commentstyle=\color{rdgrey}\itshape,
  numbers=none, xleftmargin=2pt, framexleftmargin=2pt,
  aboveskip=6pt, belowskip=6pt,
}

% Decision box: used for every choice that a reader might otherwise have to
% reverse-engineer from the code.
\newtcolorbox{decision}[1]{
  colback=rdbg, colframe=rdblue, boxrule=0.6pt, arc=2pt,
  left=6pt, right=6pt, top=5pt, bottom=5pt,
  fonttitle=\bfseries\small, title={Decision --- #1}}

\newtcolorbox{caveat}[1]{
  colback=white, colframe=rdred, boxrule=0.6pt, arc=2pt,
  left=6pt, right=6pt, top=5pt, bottom=5pt,
  fonttitle=\bfseries\small, title={#1}}

\newcommand{\unit}[1]{\,\mathrm{#1}}
\newcommand{\topicname}[1]{\texttt{\small #1}}
\newcommand{\code}[1]{\texttt{\small #1}}

% figure with a standard width and placement
\newcommand{\archfig}[3]{%
  \begin{figure}[H]\centering
  \includegraphics[width=#3\textwidth]{figures/#1.pdf}
  \caption{#2}\label{fig:#1}\end{figure}}
